% !PRINTED 41
\weberpage{41}

If we substitute $B_1, B_2 \ldots B_\mu$ for $\Theta$ here, it follows that $p_1 m$
is divisible by each of the numbers $b_1, b_2 \ldots b_\mu$. Hence $p_1$ is
divisible by $p_1'$ and $m$ by the least common
multiple of $b_1', b_2' \ldots b_\mu'$. In this argument we may now
interchange $A$ with $B$ throughout. Hence $p_1'$ must also be divisible
by $p_1$, and therefore
\weq{5}{p_1 = p_1',}
and moreover it follows that $m$ is also the least common
 multiple of $b_1', b_2' \ldots b_\mu'$.

In order to prove our theorem in general from this, we assume
that it has been proved for some $s$ that
\weq{6}{p_1 = p_1', p_2 = p_2', \ldots p_{s-1} = p_{s-1}'}
must hold. By (4), for every element $\Theta$:
\weq{7}{\Theta^{p_s m} = A_1^{\alpha_1 p_s m} A_2^{\alpha_2 p_s m} \ldots A_{s-1}^{\alpha_{s-1} p_s m},}
in which $m$, as above, is the least common multiple of
$a_1', a_2' \ldots a_\nu'$ and at the same time of $b_1', b_2' \ldots b_\mu'$, and hence is not
divisible by $p$.

We now determine the number of mutually distinct
 elements contained in the form (7).

If we let $\alpha_1$ run through the sequence of numbers $0, 1, 2 \ldots \dfrac{p_1}{p_s} - 1$,
the elements
\weq{8}{1,\ A_1^{p_s m},\ A_1^{2 p_s m} \ldots A_1^{\left(\frac{p_1}{p_s}-1\right) p_s m}}
are all distinct from one another, while $A_1^{\frac{p_1}{p_s} p_s m}$ becomes $= 1$
again, so that all the other powers $A_1^{\alpha_1 p_s m}$ are found again in the
sequence (8). We conclude in the same way with regard to the remaining
factors of (7), from which the exact number of the mutually
 distinct elements contained in $\Theta^{p_s m}$
\weq{9}{z = \frac{p_1}{p_s} \frac{p_2}{p_s} \ldots \frac{p_{s-1}}{p_s}}
is found.

Of the two numbers $p_s$, $p_s'$, if they are not equal,
one will be the larger. We shall therefore assume that
\weq{10}{p_s' \geqq p_s.}

We now express $\Theta$ in terms of the base $B$, and set
\weq{11}{\Theta^{p_s m} = B_1^{\beta_1 p_s m} B_2^{\beta_2 p_s m} \ldots B_\mu^{\beta_\mu p_s m},}
%CONT>>
